\documentclass[11pt,a4paper]{article}
\usepackage{amsmath,amsfonts,amssymb}
\usepackage[ngerman]{babel}
\usepackage[T1]{fontenc}
\usepackage[utf8]{inputenc}
\usepackage{longtable}
\usepackage{ae}
\usepackage{geometry}
\geometry{verbose,a4paper,tmargin=25mm,bmargin=25mm,lmargin=20mm,rmargin=20mm,nomarginpar}
\usepackage{tikz}
\usetikzlibrary{patterns}
\usepackage{graphicx}
\usepackage{mathrsfs}
\usepackage{array}
\usepackage{paralist}
\usepackage[pdfpagelabels]{hyperref}
\usepackage{pdfpages}
\usepackage{hyperref}
\usepackage{marvosym}
\usepackage{amssymb}
\usepackage{amsmath}
\usepackage{float}


\begin{document}
\newcommand{\angstrom}{\mbox{\normalfont\AA}}
\newcommand*{\vertbar}{\rule[1ex]{0.5pt}{2.5ex}}
\newcommand*{\horzbar}{\rule[.5ex]{2.5ex}{0.5pt}}
\newcommand{\sg}{\vspace*{0.15cm}}
\newcommand{\pic}[4]{\begin{figure}[ht] \centering \includegraphics[width=#1\textwidth]{#2}\caption{#3}\label{#4}  
\end{figure}}
\newcommand{\picoc}[3]{\begin{figure}[ht] \centering \includegraphics[width=#1\textwidth]{#2}\label{#3}  
\end{figure}}

\pagenumbering{Alph}
\begin{titlepage}

\newcommand{\HRule}{\rule{\linewidth}{0.5mm}} % Defines a new command for the horizontal lines, change thickness here

\center % Center everything on the page
 
%----------------------------------------------------------------------------------------
%	HEADING SECTIONS
%----------------------------------------------------------------------------------------

\textsc{\LARGE Albert-Ludwigs-Universität Freiburg}\\[1.5cm] % Name of your university/college
\textsc{\Large Anfänger Praktikum I}\\[0.5cm] % Major heading such as course name
\textsc{\large Versuch 17}\\[0.5cm] % Minor heading such as course title

%----------------------------------------------------------------------------------------
%	TITLE SECTION
%----------------------------------------------------------------------------------------

\HRule \\[0.4cm]
{ \huge \bfseries   Physikalisches Pendel, Trägheitsmomente und Steinerscher Satz}\\[0.4cm] % Title of your document
\HRule \\[1.5cm]

%----------------------------------------------------------------------------------------
%	IMAGE
%----------------------------------------------------------------------------------------

\includegraphics[width=200pt]{logo}\\[1cm]

 
%----------------------------------------------------------------------------------------
%	AUTHOR SECTION
%----------------------------------------------------------------------------------------
\mbox{}
\vfill
\begin{minipage}{0.4\textwidth}
\begin{flushleft} \large
\emph{Versuchsteilnehmer:}\\
\underline{Andreas \textsc{Weber}},\\ \underline{Clemens \textsc{Lauby}} % Your name
\end{flushleft}
\end{minipage}
~
\begin{minipage}{0.4\textwidth}
\begin{flushright} \large
\emph{Tutor: } \\
Vorname \textsc{Nachname} % Supervisor's Name
\end{flushright}
\end{minipage}\\[1cm]

% If you don't want a supervisor, uncomment the two lines below and remove the section above
%\Large \emph{Author:}\\
%John \textsc{Smith}\\[3cm] % Your name

%----------------------------------------------------------------------------------------
%	DATE SECTION
%----------------------------------------------------------------------------------------

{\large \today}\\ % Date, change the \today to a set date if you want to be precise

%----------------------------------------------------------------------------------------
%	LOGO SECTION, if you want no logo leave this as is
%----------------------------------------------------------------------------------------

%\includegraphics[width=200pt]{logo}\\[1cm] % Include a department/university logo - this will require the graphicx package
 
%----------------------------------------------------------------------------------------

\vfill % Fill the rest of the page with whitespace
\end{titlepage}



\newpage
\pagenumbering{Alph}
\numberwithin{equation}{section}
\tableofcontents
\newpage
\pagenumbering{arabic}
\section{Versuchsbeschreibung}
\subsection{Ziel des Versuches}
Das Ziel des Versuches ist ....
\subsection{Physikalischer Zusammenhang}
       \subsubsection{Physikalisches Pendel}
       \paragraph{Periodendauer des Physikalischen Pendels}$\\$
       
       
      Die Herleitung der Periodendauer des physikalischen Pendels erfolgt über die Lösung seiner Bewegungsgleichung. Die Bewegungsgleichung leitet sich über einen Zusammenhang des Drehmomentes her.\\
      
       Das Drehmoment ist gegeben durch:
      \begin{align}
      	\vec{M} = \vec{r}\times\vec{F} = \frac{{d\vec{l}}}{dt} \\
      	\text{mit         }        \vec{l} = I \cdot \vec{\omega} \text{  und  }  \vec{\omega} = \dot{\vec{\varphi}} \notag\\
      	\text{ergibt sich:   }  \vec{M} = I \cdot \ddot{\vec\varphi}\notag
      	  \end{align}
Bzgl. eines Starren Körpers in Form eine Stabes mit Trägheitsmoment $I_a$, Masse $m_s$ und $d$ als Abstand zwischen Aufhängepubnkt und Schwerpunkt, ist der Betrag des Drehmoments folglich
\begin{equation}
		M =- I_a \ddot{\varphi}
\end{equation}
        Aus (1) ergibt sich:
\begin{align}
 		\vec{M}= \vec{F_G} \times \vec{d} \overset{\text{(Betrag)}}{\underset{\text{ }}{\rightarrow}} M = m_sgd \sin\varphi
\end{align}
Gleichsetzen von (2) und (3) ergibt
\begin{align*}
		 m_sgd \sin\varphi&= - I_a \ddot{\varphi} \\
 		\Leftrightarrow I_a \ddot{\varphi} + m_sgd \sin\varphi&= 0
\end{align*}
Für kleine Winkel $\varphi$ folgt
\begin{align*}
		I_a \ddot{\varphi} + m_sgd \varphi&= 0\\
		 \Leftrightarrow \ddot\varphi + \left(\frac{m_sgd}{I_a}\right)\varphi &= 0
	\end{align*}
Diese Differentialgleichung hat die Form einer allgemeinen Schwingungsgleichung
\begin{align*}
	\ddot{x}+\omega^2 x = 0
\end{align*}
mit $\omega^2 = \left(\frac{m_sgd}{I_a}\right)$. Die Lösung dieser Differentialgleichung ist eine harmonische Schwingung.\\
Die Kreisfrequenz $\omega$ liefert uns durch folgende Zusammenhänge die Periodendauer
\begin{equation*}
	T = \frac{1}{f} = \frac{1}{2\pi}\sqrt{\frac{I_a}{m_sgd}} \hspace{2cm} \text{mit}\hspace{0.5cm}   f=\frac{\omega}{2\pi}
\end{equation*}




	
       \subsubsection{Drehpendel}
       
       
       
       
       
       
       
       
       
       
       
       
       
       
       
       
       

\subsection{Versuch zum Aufgabenteil 3}
      \subsubsection{Versuchsaufbau}
      \subsubsection{Durchführung}

\subsection{Versuch zum Aufgabenteil 4}
     \subsubsection{Versuchsaufbau}
     \subsubsection{Durchführung}
\section{Auswertung und Fehleranalyse}
	\subsection{Physikalisches Pendel}
	
	Im ersten Teil unseres Versuches haben wir 30 mal die Periodendauer einer Schwingung des Pendels
	gemessen. Aus den 30 Messwerten errechnen wir den Mittelwert mit:
	
	\begin{equation}
	\bar{T} = \frac{1}{n} \sum_{i=1}^{n} T_i		
	\end{equation}
	Die Standardabweichung mit:
	\begin{equation}
	s_T=\sqrt{\frac{1}{n-1}\sum_{i=1}{n} (T_i - \bar{T_s})^2}
	\end{equation}
	Die Standardabweichung des Mittelwerts mit:
	\begin{equation}
	s_{\bar{T}}=\frac{s_T}{\sqrt{n}}
	\end{equation}
	Das Ergebnis bezüglich unserer 30 Messwerte (siehe Anhang) ist also:
	\begin{equation*}
	\bar{T_e}=(1.5850 \pm 0.0081)s
	\end{equation*}
	\\
	Im zweiten Teil des Versuches haben wir die Schwingzeit von 30 Perioden gemessen. \\
	Dabei haben wir unsere Messungenauigkeit auf $\pm 0.1$s abgeschätzt. Also ist $T_{30}=(48.03\pm 0.1)s$
	\\
	Die Schwingungsdauer einer Periode berechnet sich als:
	\begin{equation*}
	\bar{T}=\frac{T_{30}}{30}
	\end{equation*}
	Mit einer Standardabweichung von $$\frac{s_{T_{30}}}{30}=\frac{0.1}{30}[s]\approx 0.003[s]$$
	Also ist $\bar{T_r}=(1.601\pm 0.003)s$. \\
	Um die Verträglichkeit der beiden Periodendauern zu bestimmen verwenden wir:
	\begin{equation}
	t=\frac{ |x_0 -y_0|}{\sqrt{(\Delta x)^2 + (\Delta y)^2}}
	\end{equation}
	Für unsere Werte und Unsicherheiten ergibt sich $t=1.96<2$. Also sind unsere Ergebnisse verträglich.
	\newpage 
	Vergleicht man das Ergebnis der 30 Einzelmessungen mit dem der Messung von 30 Schwingungen,
	so fällt auf, dass die Standardabweichung bei der Messung von einer einzelnen Periode deutlich 
	größer ist als die Standardabweichung einer Periode bei der Messung von 30 Schwingungen.
	Bezüglich diesem Aspekt ist die Messung mehrerer Perioden der Messung einer einzelnen Periode überlegen.
	\\
	Das Ergebnis des durchgeführten T-Tests ($t=1.96$) ist zwar noch innerhalb der Akzeptanz, aber
	ziemlich nah an dem kritischen Wert von 2.
	Möglicher Grund dafür könnte eine zu niedrige Einschätzung unsere statistischer Messungenauigkeit bei der Messung von 30 Schwingungen sein.
	\\Das mehrfache durchführen der Messung von 30 Perioden hätte den Vorteil, dass man die Messungenauigkeit berechnen könnte und sich nicht auf eine Abschätzung verlassen müsste.
	\\
	Aus diesem Grund wäre eine Mischung der beiden, von uns durchgeführten, Messmethoden ideal um ein
	besonders genaues Ergebnis zu erhalten.
	
	\subsubsection{Bestimmung der Periodendauer}
	Aus der physikalischen Herleitung wissen wir, das T mit:
	\begin{equation}
	T=2\pi \sqrt{\frac{2l}{3g}}
	\end{equation}
	berechnet werden kann.\\
	Mit $l=(0.954\pm 0.001)$ (siehe Protokollheft) \\ und $g=9.81$(Demtröder: Experimentalphysik 1,Springer,2018,s.25)
	sowie der Gaußschen Fehlerfortpflanzung:
	\begin{equation}
	\Delta z = \sqrt{(\frac{\partial f}{\partial x}\Delta x)^2 +(\frac {\partial f}{\partial y} \Delta y)^2+...} \text{   für   } z=f(x,y,...)
	\end{equation}
	ergibt sich T rechnerisch zu: $T_{rech}=(1.599830\pm0.000054)$s\\
	{\bf Diskussion der Abweichung zwischen berechneter und gemessener Schwingungszeit}\\
	Ein T-Test (2.4) zwischen den Werten $T_e$ und $T_{rech}$ ergibt $t=1.82<2$(wobei $\Delta T_{rech}$ im Vergleich zu $\Delta T_e$ sehr klein ist und deswegen in der Berechnung von t vernachlässigt werden kann).\\
	Also sind die Abweichungen zwischen dem berechnetem Wert und dem experimentell bestimmten Wert
	verträglich miteinander. Unser , durch Messen bestimmtes, Ergebnis stimmt also mit der theoretischen Herleitung überein.
	
	
	\subsection{Drehpendel}
	
	
	
	
	
	
	
	
	
	
\subsection{Messwerte}



\newpage
\section{Anhang}
%\bibliographystyle{apacite}

\nocite{*}
\newpage
\section{Literatur}
\newpage
%\includepdf[pages={1-}]{abc.pdf}



\end{document}

